\section{Удиртгал}
\label{sec:introduction}

Багц сонголтын дундаж--вариацын онолын суурийг \citet{markowitz1952} тавьсан бөгөөд багцын эрсдэлийг тодорхойлоход хөрөнгийн өгөөжийн ковариацын матриц гол үүрэгтэй. Глобал хамгийн бага вариацтай багц (Global Minimum Variance Portfolio, GMVP) нь хүлээгдэж буй өгөөжийн векторыг үнэлэхгүйгээр багцын вариацыг хамгийн бага байлгах жинг тодорхойлдог. Гэвч эх олонлогийн ковариацын матриц үл мэдэгдэх тул түүврээс үнэлэх шаардлагатай. Түүврийн ковариацын матриц $N$ хөрөнгийн хувьд $N(N+1)/2$ ялгаатай параметртэй бөгөөд $N/T$ харьцаа өсөхөд үнэлгээний алдаа, матрицын нөхцөлт тоо $(\lambda_{max}/\lambda_{min})$ болон багцын жингийн тогтворгүй байдал зэрэгт хүндрэл үүсгэж болно \citep{ledoitwolf2004a}. Ковариацын матрицын урвуу олдох боломжтой, $N/T$ бага байсан ч ковариацын үнэлгээний алдаа GMVP-ийн түүврээс гадуурх гүйцэтгэлд нөлөөлөх боломжтой. \citet{demiguel2009} сарын өгөөжийн долоон өгөгдлийн сан дээр багц сонголтын 14 загварыг $1/N$ жишигтэй харьцуулж, оновчлолын аргууд уг жишгээс бүх нөхцөлд тогтвортой давуу байгаагүйг харуулсан.

Ковариацын матрицын үнэлгээний алдааг бууруулах түгээмэл арга нь агшилтын үнэлгээ юм. Шугаман агшилт түүврийн ковариацын матриц $S$-ийг бүтэцтэй зорилтот матрицтай хослуулдаг. Энэ чиглэлийн судалгаанд нэг хүчин зүйлийн загварт суурилсан зорилтот матриц \citep{ledoitwolf2003}, тогтмол вариацын зорилтот матриц \citep{ledoitwolf2004a}, тогтмол корреляцын зорилтот матриц \citep{ledoitwolf2004b}, мөн диагональ зэрэг бүтцийн хувилбарууд \citep{SchaferStrimmer2005,AndoXiao2023} багтана. Эдгээр нь хөрөнгүүдийн дисперс, ковариац болон нийтлэг хүчин зүйлийн талаарх ялгаатай таамаглал ашигладаг. Харин шугаман бус агшилт нь ковариацын матрицын түүврийн хувийн утгуудыг ялгаатай хэмжээгээр засварлах боломжтой \citep{ledoitwolf2012,ledoitwolf2020}. Иймээс аль үнэлэгч илүү сайн ажиллах нь зөвхөн агшилт хэрэглэсэн эсэхээс бус, зорилтот бүтцийн тохироо, өгөгдлийн шинж чанар болон багцын оновчлолын нөхцөлөөс хамаарч болно \citep{ledoitwolf2022}.

Ковариацын үнэлэгчийн статистикийн алдагдлыг бууруулах нь багцын түүврээс гадуурх эрсдэлийг заавал бууруулна гэсэн үг биш. \citet{ledoitwolf2017} багцын оновчлолд зориулсан шугаман бус агшилтын аргыг судалж, түүврээс гадуурх гүйцэтгэлийг үнэлсэн. \citet{shi2025} сарын өгөөж болон 120 сарын шилжих цонх ашиглан агшилтын эрчмийг cross-validation-д суурилсан negative log-likelihood шалгуураар сонгох аргыг санал болгож, харьцангуй олон хөрөнгөтэй, богино борлуулалтыг зөвшөөрсөн GMVP-ийн хүрээнд диагональ болон нэг хүчин зүйлийн зорилтот матрицын давуу талыг тайлагнасан. Эдгээр судалгаа нь ковариацын үнэлгээний зорилго, агшилтын эрчим болон багцын бодит гүйцэтгэлийн хоорондын холбоог тодруулсан боловч тэдгээрийн үр дүнг бага $N/T$-тэй, богино борлуулалт хориглосон багцад шууд шилжүүлэн үзэх боломжгүй.

Багцын жингийн хязгаарлалт өөрөө үнэлгээний алдааны нөлөөг бууруулж болох нь энэ асуудлыг улам чухал болгодог. \citet{JagannathanMa2003} жинг сөрөг бус байлгах хязгаарлалт нь ковариацын үнэлгээнд агшилттай төстэй нөлөө үзүүлж болохыг тайлбарлаж, ийм нөхцөлд түүврийн ковариацад суурилсан GMVP фактор болон агшилтын үнэлэгчдэд суурилсан багцуудтай ойролцоо гүйцэтгэлтэй байж болохыг харуулсан. \citet{ledoitwolf2017} мөн багцын хязгаарлалттай болон хязгаарлалтгүй тохиолдлын үр дүнг ялган авч үзсэн. Сүүлийн үеийн \citet{DeNardKostovic2026} судалгаа агшилтын эрчмийг өгөгдөлд суурилан суралцах аргыг санал болгож, unrestricted болон long-only GMVP-ийг харьцуулсан бөгөөд жингийн хязгаарлалт нь нэмэлт агшилтын үр ашгийг багасгаж болохыг тэмдэглэсэн. Гэвч тус судалгааны үндсэн эмпирик дизайн өдрийн өгөөжид тулгуурласан тул сарын өгөөжтэй, бага $N/T$ нөхцөлд адил дүгнэлт гарах эсэх нь тусдаа эмпирик асуулт юм.

Эдгээр судалгаанаас харахад ковариацын агшилтын онол, өндөр хэмжээст багцын хэрэглээ болон жингийн хязгаарлалтын нөлөөг өргөн судалжээ. Харин сарын өгөөжид суурилсан, $N/T$ харьцаа бага, long-only GMVP-ийн хүрээнд агшилтын өөр өөр зорилтот бүтэцтэй үнэлэгчдийг хөрөнгүүдийн корреляцийн ялгаатай түвшин болон олон багцын бүрэлдэхүүнд зэрэгцүүлэн шалгах нь ач холбогдолтой хэвээр байна. Ялангуяа нэг багцад илэрсэн дисперсийн бууралт бусад багцад давтагдах эсэх, мөн ажиглагдсан ялгаа статистикийн тодорхой бус байдлыг харгалзсаны дараа хадгалагдах эсэхийг тусад нь үнэлэх шаардлагатай. Энэхүү судалгаа нь өмнөх ажлууд агшилтыг эсвэл long-only GMVP-ийг огт судлаагүй гэж үзэхгүй; харин эдгээр нөхцөлийг нэг эмпирик хүрээнд хослуулан харьцуулж байгаагаараа ялгаатай.

Судалгааны зорилго нь түүврийн ковариацын матриц (Sample), тогтмол вариацын (Var), диагональ (Diag), тогтмол корреляцын (Corr), нэг хүчин зүйлийн (SF) зорилтот матрицад суурилсан дөрвөн шугаман агшилтын үнэлэгч болон аналитик шугаман бус агшилтын үнэлэгчийг (NLS) богино борлуулалтыг хориглосон GMVP-ийн ижил нөхцөлд харьцуулах явдал юм. Сарын өгөөжид үндэслэн тус бүр 10 хувьцаатай, дундаж корреляци өндөр, дунд, бага гурван багц байгуулж, 120 сарын шилжих сургалтын цонхоор 2022 оны 6 дугаар сараас 2026 оны 9 дүгээр сар хүртэлх 52 сарын түүврээс гадуурх гүйцэтгэлийг үнэлнэ. Үндсэн шалгуур нь түүврээс гадуурх хэлбэлзэл болон Sample аргатай харьцуулсан дисперсийн бууралт (RVR) бөгөөд Sharpe-ийн харьцаа, багцын жингийн эргэлт (turnover), төвлөрлийг мөн авч үзнэ.

Судалгааны эмпирик хувь нэмэр нь дараах харилцан нөхөх шинжилгээнд оршино. Нэгдүгээрт, корреляцийн гурван ангилалд үнэлэгчдийн харьцангуй гүйцэтгэлийг харьцуулж, сургалтын цонхыг 60, 90, 120 сар болгон өөрчилснөөр $N/T$-ийн мэдрэмжийг шалгана. Хоёрдугаарт, ангилал тус бүрээс 50, нийт 150 нэмэлт багц сонгож, үр дүн багцын бүрэлдэхүүнээс хамаарах эсэхийг үнэлнэ. Гуравдугаарт, HAC болон circular moving-block bootstrap аргаар статистикийн тодорхой бус байдлыг тооцож, олон харьцуулалтын $p$-утгад Holm-ийн засвар хэрэглэнэ. Дөрөвдүгээрт, $1/N$ жишиг багц болон арилжааны зардлын дараах гүйцэтгэлийг нэмэлтээр харьцуулна.

Үндсэн гурван багцын үр дүнгээр өндөр корреляцитай багцад Var, дунд корреляцитай багцад Corr үнэлэгч Sample аргатай харьцуулахад түүврээс гадуурх дисперсийг тус тус 5.81\% болон 2.26\%-иар бууруулсан бол бага корреляцитай багцад агшилтын таван үнэлэгчийн RVR бүгд сөрөг байв. Нэмэлт 150 багцын үр дүн нь агшилтын харьцангуй гүйцэтгэл багцын бүрэлдэхүүнээс хамаарч өөрчлөгдөж болохыг харуулсан. Олон харьцуулалтын засварын дараа агшилтын үнэлэгчдийн Sample аргаас, мөн үндсэн багцуудын GMVP стратегиудын $1/N$ жишгээс статистикийн ач холбогдолтой давуу байдал тогтоогдоогүй. Иймээс ажиглагдсан эрсдэлийн ялгааг бүх багцад хамаарах баталгаатай давуу тал бус, сонгосон өгөгдөл болон багцын бүтцийн хүрээнд илэрсэн эмпирик үр дүн гэж тайлбарлана.
